\documentclass[../../../../main.tex]{subfiles}
\begin{document}

\begin{flushright}
\footnotesize\textit{【自動文字起こし・要確認】}
\end{flushright}

\noindent
[\textbf{解}] 3辺の長さ $a, b, c$ ($a \leqq b \leqq c$) とすると, 三角不等式から
\[
a \leqq b \leqq c < a + b
\]
が成り立つ.
\[
a + b - c \geqq 20 + 20 - 36 \geqq 0
\]
だから, 問題文の中にある長さのうち任意に3つ持ってくると, 三角形をつくることができる. から, 9つの辺の長さから3つをえらぶ方法を考えて

\begin{itemize}
    \item[$1^\circ$] 3辺とも違う長さ $\dots {}_9\mathrm{C}_3 = 84$ 通り
    \item[$2^\circ$] 2等辺三角形 $\dots 2 \cdot {}_9\mathrm{C}_2 = 72$ 通り
    \item[$3^\circ$] 正三角形 $\dots 9$ 通り
\end{itemize}

よって
\[
84 + 72 + 9 = 165 \text{ 通り}
\]

\paragraph*{[本番のミス]}
“相似”と“合同”をまちがえていた

\newpage

\end{document}
